%
\usepackage{color}
\usepackage{iftex}
\ifPDFTeX
  \usepackage[sfdefault]{carlito}
  \renewcommand{\familydefault}{\sfdefault}
\else
  % Ana yazı tipi inonutez.cls içinde ayarlanır. Burada tekrar
  % Calibri/Consolas araması yapılmaz; aksi halde Calibri bulunmayan
  % sistemlerde XeLaTeX gereksiz font hatası verebilir.
  \IfFontExistsTF{consola.ttf}{\setmonofont{consola.ttf}[BoldFont=consolab.ttf,ItalicFont=consolai.ttf,BoldItalicFont=consolaz.ttf,Scale=MatchLowercase]}{
    \IfFontExistsTF{DejaVuSansMono.ttf}{\setmonofont{DejaVuSansMono.ttf}[Scale=MatchLowercase]}{}
  }
\fi
\usepackage{amssymb,amsmath,amsbsy,bm}
\usepackage{caption}            
\usepackage{graphics}
\usepackage{wrapfig}
\usepackage{epsfig}
\usepackage{enumerate}
\usepackage{rotating}
\usepackage{multirow}
\usepackage{colortbl}
\usepackage{latexsym}  
\usepackage{rotating}

% ----------------------------------------------------------------
% Tez genel preamble ayarlari. Ana tez.tex dosyasini sade tutmak
% icin paketler, kaynakca destekleri ve dizin ayarlari burada.
% ----------------------------------------------------------------
\usepackage{afterpage}
\usepackage{tocloft}
\usepackage{graphicx}

\makeatletter
\renewcommand{\cftfigpagefont}{\bfseries}
\renewcommand{\cfttabpagefont}{\bfseries}
\renewcommand{\cftpartleader}{\cftdotfill{\cftdotsep}}
\renewcommand{\cftchapleader}{\cftdotfill{\cftdotsep}}

\setlength{\cftfigindent}{0pt}
\renewcommand{\cftfigpresnum}{\bfseries Şekil }
\setlength{\cftfignumwidth}{5.5em}
\renewcommand{\cftfigaftersnum}{. }
\renewcommand{\cftdotsep}{1}

\setlength{\cfttabindent}{0pt}
\renewcommand{\cfttabpresnum}{\bfseries Tablo }
\setlength{\cfttabnumwidth}{5em}
\renewcommand{\cfttabaftersnum}{. }

\setlength{\cftchapnumwidth}{1.6em}
\setlength{\cftsecnumwidth}{1.6em}
\setlength{\cftsubsecnumwidth}{2.4em}
\setlength{\cftsubsubsecnumwidth}{3.2em}

\setlength{\cftpartindent}{0em}
\setlength{\cftchapindent}{0em}
\setlength{\cftsecindent}{1.1em}
\setlength{\cftsubsecindent}{1.9em}
\setlength{\cftsubsubsecindent}{2.7em}

\setlength{\cftbeforechapskip}{1.055em}
\setlength{\cftbeforesecskip}{0.035em}
\setlength{\cftbeforesubsecskip}{0.035em}

\setlength{\cfttabindent}{0em}
\renewcommand{\cfttabpresnum}{\bfseries Tablo }
\setlength{\cfttabnumwidth}{5.5em}
\renewcommand{\cfttabaftersnum}{. }

\newcommand{\dizinBaslikYazisi}{\sffamily\bfseries\fontsize{14}{16.8}\selectfont}
\renewcommand{\contentsname}{\makebox[\textwidth][c]{\dizinBaslikYazisi \if@Ingilizce TABLE OF CONTENTS\else İÇİNDEKİLER\fi}}
\renewcommand{\listfigurename}{\makebox[\textwidth][c]{\dizinBaslikYazisi \if@Ingilizce LIST OF FIGURES\else ŞEKİLLER DİZİNİ\fi}}
\renewcommand{\listtablename}{\makebox[\textwidth][c]{\dizinBaslikYazisi \if@Ingilizce LIST OF TABLES\else TABLOLAR DİZİNİ\fi}}
\newcommand{\dizinSayfaBasligi}{%
  \par\nobreak
  \vspace{0.5\baselineskip}%
  \noindent\makebox[\textwidth][r]{\bfseries\underline{\if@Ingilizce Page\else Sayfa\fi}}%
  \par\nobreak
}
\renewcommand\cfttoctitlefont{}
\renewcommand\cftloftitlefont{}
\renewcommand\cftlottitlefont{}
\renewcommand\cftaftertoctitle{\dizinSayfaBasligi}
\renewcommand\cftafterloftitle{\dizinSayfaBasligi}
\renewcommand\cftafterlottitle{\dizinSayfaBasligi}
\renewcommand{\cftaftertoctitleskip}{10pt}
\renewcommand{\cftafterloftitleskip}{10pt}
\renewcommand{\cftafterlottitleskip}{10pt}
\makeatother

% Kaynak stili documentclass icindeki apa veya num secenegi ile belirlenir.
\makeatletter
\if@APAStyle
  \usepackage[style=apa, backend=biber, sortcites=true, maxnames=20, minnames=1, dashed=false]{biblatex}
  \AtBeginBibliography{%
    \oneandonehalf
    \sffamily\normalsize\selectfont
    \setlength{\bibhang}{1.25cm}%
    \renewcommand*{\mkbibnamefamily}[1]{#1}%
    \renewcommand*{\mkbibnamegiven}[1]{#1}%
  }
  \addbibresource{kaynaklar.bib}
  \defbibheading{bibliography}{%
    \makeatletter
    \chapter*{\if@Ingilizce REFERENCES\else KAYNAKLAR\fi}
    \addtocontents{toc}{\protect\addvspace{-10pt}}
    \addcontentsline{toc}{chapter}{\bf{\if@Ingilizce REFERENCES\else KAYNAKLAR\fi}}
    \makeatother
    \sffamily\normalsize\selectfont
  }
\else
  \providecommand{\parencite}[1]{\cite{#1}}
  \providecommand{\textcite}[1]{\cite{#1}}
\fi
\makeatother

% BEGIN TEZ_GUI_THEOREM_ENVIRONMENTS
\usepackage{amsthm}
\makeatletter
\renewcommand{\proofname}{\if@Ingilizce Proof\else İspat\fi}
\makeatother
\theoremstyle{plain}
\newtheorem{theorem}{Teorem}[section]
\newtheorem{lemma}[theorem]{Lemma}
\newtheorem{proposition}[theorem]{Önerme}
\newtheorem{corollary}[theorem]{Sonuç}
\theoremstyle{definition}
\newtheorem{definition}[theorem]{Tanım}
\newtheorem{example}[theorem]{Örnek}
\newtheorem{problem}[theorem]{Problem}
\theoremstyle{remark}
\newtheorem{remark}[theorem]{Uyarı}
% END TEZ_GUI_THEOREM_ENVIRONMENTS

\def\be{\begin{equation}} %
\def\ee{\end{equation}}%
\def\beq{\begin{eqnarray}}%
\def\eeq{\end{eqnarray}}%
\def\bse{\begin{subequations}}%
\def\ese{\end{subequations}}%
\def\nonu{\nonumber}%
\def\psibar{\overline{\psi}} %
\def\Delslash{\partial\!\!\!\!\!/}%
\def\Gslash{G\!\!\!\!\!/}%
\def\Lmbdslash{\lambda\!\!\!\!\!/}%
\def\Jslash{J\!\!\!\!\!/}%
\def\pslash{p\!\!\!\!/}%
\def\qslash{q\!\!\!\!/}%
\def\kslash{k\!\!\!\!/}%
\def\lb{\left[}
\def\rb{\right]}
\def\lp{\left(}
\def\rp{\right)}

\def\Lag{{\cal L}}    % 
\def\D{{\cal D}}      % 
\def\H{{\cal H}}      % 
\def\Z{{\cal Z}}      % 
\def\S{{\textbf{S}}}  % 
\def\d{{\rm d}}%
\def\dt{{\rm{dt}}}%
\def\Tr{{\rm{Tr}}}%
\def\gm{\gamma^{\mu}}
\def\ga{\gamma^{\alpha}}
\def\gb{\gamma^{\beta}}
\def\gn{\gamma^{\nu}}
\def\gs{\gamma^{\sigma}}
\def\gl{\gamma^{\lambda}}
\def\gr{\gamma^{\rho}}
\def\gd{\gamma^{\delta}}
